\documentclass[10pt]{article}
\usepackage[T1]{fontenc}
\usepackage[utf8]{inputenc}
\usepackage[french]{babel}
\usepackage{amsmath,amssymb,mathrsfs,mathpazo}
\usepackage{geometry}
\geometry{a4paper,margin=2cm}
\usepackage{graphicx,array,tabularx,tabularray,colortbl,multirow,multicol,booktabs,enumitem,pifont,ulem,fancybox,diagbox,cancel}
\usepackage{pstricks,pstricks-add,pst-plot,pst-tree,pst-node,pst-text,pst-func,pst-math,pst-bspline,pst-3dplot,pst-eucl,pst-solides3d}
\let\pscalculate\undefined
\usepackage{pst-bezier}
\usepackage{tikz,pgfplots,tkz-tab}
\usetikzlibrary{arrows,arrows.meta,patterns,calc,positioning,trees,decorations.markings,shapes.geometric,shapes.misc,shapes.arrows,intersections,decorations.pathmorphing,decorations.pathreplacing,backgrounds,shadows,babel}
\pgfplotsset{compat=1.18}
\usepgfplotslibrary{fillbetween}
\usepackage[np]{numprint}
\usepackage{listings,makecell,pas-tableur,twoopt}
\let\barre\undefined
\usepackage[pstricks]{tablvar}
\definecolor{darkgreen}{RGB}{0,100,0}
\providecommand{\red}{\color{red}}
\providecommand{\blue}{\color{blue}}
\providecommand{\green}{\color{green}}
\providecommand{\black}{\color{black}}
\providecommand{\white}{\color{white}}
\providecommand{\gray}{\color{gray}}
\providecommand{\cyan}{\color{cyan}}
\providecommand{\magenta}{\color{magenta}}
\providecommand{\yellow}{\color{yellow}}
\providecommand{\violet}{\color{violet}}
\providecommand{\orange}{\color{orange}}
\providecommand{\euro}{\texteuro}
\providecommand{\EUR}{\texteuro}
\providecommand{\ieme}{\textsuperscript{e}}
\providecommand{\degres}{\ensuremath{^\circ}}
\providecommand{\degree}{\ensuremath{^\circ}}
\providecommand{\vv}[1]{\overrightarrow{#1}}
\providecommand{\Coord}[1]{#1}
\providecommand{\decosix}{}
\providecommand{\starredbullet}{\ensuremath{\bullet}}
\pagestyle{empty}
\setlength{\parindent}{0pt}
\usepackage{adjustbox}\DeclareUnicodeCharacter{2236}{:}\DeclareUnicodeCharacter{2190}{\ensuremath{\leftarrow}}
\providecommand{\R}{\mathbb{R}}
\providecommand{\N}{\mathbb{N}}
\providecommand{\Z}{\mathbb{Z}}
\providecommand{\Q}{\mathbb{Q}}
\providecommand{\D}{\mathbb{D}}
\providecommand{\C}{\mathbb{C}}
\providecommand{\e}{{\mathrm{e}}}
\providecommand{\vect}[1]{\overrightarrow{#1}}
\providecommand{\vectt}[1]{\overrightarrow{\mathrm{#1}}}
\providecommand{\vvt}[1]{\overrightarrow{\mathrm{#1}}}
\providecommand{\np}[2][]{{#2~#1}}
\providecommand{\vv}[1]{{\overrightarrow{#1}}}
\renewcommand{\d}{{\mathrm{d}}}
\providecommand{\degres}{{^\circ}}
\providecommand{\degre}{{^\circ}}
\providecommand{\blue}{\color{blue}}
\providecommand{\red}{\color{red}}
\providecommand{\enskip}{\quad}
\providecommand{\llbracket}{\lbrack\!\lbrack}
\providecommand{\rrbracket}{\rbrack\!\rbrack}
\providecommand{\euro}{{\text{€}}}
\providecommand{\no}{{\text{n°}}}
\providecommand{\interval}[2]{{[#1;#2]}}
\providecommand{\intervFO}[2]{{[#1;#2[}}
\providecommand{\dotfill}{\ldots}

\providecommand{\ds}{\displaystyle}
\providecommand{\textding}[1]{#1}
\providecommand{\Oij}{\(\mathrm{O};\vec{\imath},\vec{\jmath}\)}
\providecommand{\Oijk}{\(\mathrm{O};\vec{\imath},\vec{\jmath},\vec{k}\)}
\providecommand{\Ouv}{\(\mathrm{O};\vec{u},\vec{v}\)}
\providecommand{\intervFF}[2]{\ensuremath{[#1;#2]}}
\renewcommand{\euro}{\ifmmode\text{\texteuro}\else\texteuro\fi}
\DeclareUnicodeCharacter{00B0}{\ifmmode^{\circ}\else\textdegree\fi}
\DeclareUnicodeCharacter{20AC}{\ifmmode\text{\texteuro}\else\texteuro\fi}
% Énoncés : Éducation nationale. Transcription et corrigés : APMEP, Denis Vergès et contributeurs. Découpage et classement : CoopMaths, Sébastien Lozano, Éric Elter et contributeurs. Adaptation de la présentation : Maths Explorer.
% CC BY-SA : https://coopmaths.fr/www/about/
% https://coopmaths.fr/alea/static/bac/2022/tex/bac_2022_10_sujet1_caledonie_2_cor.tex
\begin{document}
\noindent\textit{APMEP / CoopMaths. CC BY-SA. Présentation adaptée par Maths Explorer.}\par\medskip
\medskip









\[f(x) = x^3\e^x.\]



\medskip

\begin{enumerate}
\item On définit la suite $\left(u_n\right)$ par $u_0 = - 1$ et, pour tout entier naturel $n,\: u_{n+1} = f\left(u_n\right)$.
	\begin{enumerate}
		\item 
$\bullet~~$$u_1 = (- 1)^3 \e^{-1} = - \e^{-1} = - \dfrac{1}{\e} \approx - 0,368 $ ;

$\bullet~~$$u_2 = f\left(u_1\right) = \left(- \e^{-1}\right)^3\left(\e^{-\e^{-1}}\right) = - \e^{-3} \times \left(\e^{-\e^{-1}}\right) = - \e^{- 3 - \e^{-1}} \approx - 0,034$.
		

		\item 

\smallskip



















On a \texttt{fonc(2)}$ = u_2 \approx - 0,034$.
	\end{enumerate}
\item 
	\begin{enumerate}
		\item 
En dérivant $f(x)$ comme un produit on obtient :
		
$f'(x) = 3x^2 \e^x + x^3\e^x = x^2\e^x(3 + x)$.
		\item ~


\begin{center}
\psset{unit=1cm,arrowsize=2pt 3}
\begin{pspicture}(7,2.5)
\psframe(7,2.5)
\psline(0,2)(7,2)\psline(1,0)(1,2.5)
\uput[u](0.5,1.9){$x$} \uput[u](1.2,1.9){$- \infty$} \uput[u](4,1.9){$-3$} \uput[u](6.5,1.9){$+ \infty$} 
\rput(0.5,1){$f$}\uput[d](1.2,2){0}\uput[u](4,0){$- 27\e^{-3}$}\uput[d](6.5,2){$+ \infty$}
\psline{->}(1.5,1.5)(3.5,0.5)\psline{->}(4.5,0.52)(6.5,1.5)
\end{pspicture}
\end{center}

$\bullet~~$Quel que soit le réel $x$, \: $x^2 \geqslant 0$ et $\e^x > 0$, donc le signe de $f'(x)$ est celui de $3 + x$ qui s'annule pour $x = - 3$, d'où les deux intervalles de variations ;

$3 + x < 0 \iff x <- 3$ : sur $]- \infty~;~-3[,\: f'(x) < 0$ : la fonction $f$ est donc décroissante sur $]- \infty~;~- 3[$  ;

$3 + x > 0 \iff x >- 3$ : sur $]- 3~;~+\infty[,\: f'(x) > 0$ : la fonction $f$ est donc croissante sur $]- 3~;~+\infty[$ ;

$f(- 3) = (- 3)^3 \times \e^{-3} = - 27\e^{-3} = - \dfrac{27}{\e^3} \approx - 1,344$ est le minimum de la fonction sur $\R$.

$\bullet~~$ On a $f(0) = 0^3 \times \text{e}^0 = 0 \times 1 = 0$ ;

$\bullet~~$ De $\displaystyle\lim_{x \to + \infty} x^3 = + \infty$ et $\displaystyle\lim_{x \to + \infty} \e^{x} = + \infty$, on a par produit de limites

$\displaystyle\lim_{x \to + \infty} f(x) = + \infty$.
		\item 
		

\emph{Initialisation} : avec $u_0 = - 1$ et $u_1 \approx - 0,368$, on a bien : $ - 1 \leqslant u_0 \leqslant u_1  \leqslant 0$ : l'encadrement est vrai au rang $0$.

\emph{Hérédité} : soit $n \in \N$ et supposons que 	$- 1 \leqslant u_n \leqslant u_{n+1} \leqslant 0$.

On a vu que sur l'intervalle $]- 3~;~+\infty[$, donc a fortiori sur l'intervalle $]- 1~;~+ \infty[$, la fonction est strictement croissante.

On a donc par croissance de $f$ : $f(- 1) \leqslant f\left(u_n\right) \leqslant f\left(u_{n+1}\right)\leqslant f(0)$, ou encore :

$u_1\leqslant u_{n+1} \leqslant u_{n+2} \leqslant f(0)$	

et comme $u_1 \approx -0,338, \: \: - 1 \leqslant u_1$ et $f(0) = 0$, on a bien :

$- 1 \leqslant u_{n+1} \leqslant u_{n+2} \leqslant 0$.

La relation est vraie au rang $n + 1$.

Conclusion : l'encadrement est vrai au rang $0$ et s'il est vrai au rang $n,\:\:n \in \N$, il est encore vrai au rang $n + 1$ : d'après le principe de récurrence, pour tout entier naturel $n$, on a :
$1 \leqslant u_n \leqslant u_{n+1} \leqslant 0$.
		\item 
		
La question précédente montre que :
		
$\bullet~~$la suite $\left(u_n\right)$ est croissante ;

$\bullet~~$la suite $\left(u_n\right)$ est majorée par $0$ ;

La suite $\left(u_n\right)$ est donc convergente vers une limite $\ell$, avec $\ell \leqslant 0$.
		\item 




On résout dans $]- 1~;~0[$, (car d'après la question précédente tous les termes de la suite appartiennent à cet intervalle) l'équation  :

$f(x) = x \iff x^3\e^x = x \iff x^3\e^x - x = 0 \iff x\left(x^2\e^x - 1 \right) = 0 \iff x = 0$, car on admet que l'équation $x^2\e^x - 1 = 0$ n'a pas de solution dans l'intervalle $]- 1~;~0[$.

Conclusion $\ell = \displaystyle\lim_{n \to + \infty} u_n = 0$.
	\end{enumerate}
\end{enumerate}

\bigskip

\end{document}
